\section{Multi-dimensional distribution functions}

If $\mathbf{X}=(X_1,\ldots,X_d)$ is a $d$-dimensional
random vector, we define its $\textbf{distribution}$
to be the probability measure

\[
\mu_{\mathbf{X}}(A) = \mathrm{P}(X^{-1}(A)) = \mathrm{P}(\omega \in \Omega : X(\omega) \in A),
\qquad A \in {\cal B}(\R^d),
\]

\noindent similarly to the one-dimensional case. The measure
$\mu_X$ is also called the
$\textbf{joint distribution}$ (or
$\textbf{joint law}$) of the random variables
$X_1,\ldots,X_d$.

Once again, to avoid having to work with measures, we
introduce the concept of a $\textbf{$d
$-dimensional distribution function}$.

\begin{defn}
The $d$-dimensional distribution function of a
$d$-dimensional random vector
$\mathbf{X}=(X_1,\ldots,X_d)$ (also called the joint
distribution function of $X_1,\ldots,X_d$) is the
function $F_{\mathbf{X}}:\R^d\to[0,1]$ defined by

\begin{align*}
F_{\mathbf{X}}(x_1,x_2,\ldots,x_d)
&= \mathrm{P}( X_1 \le x_1, X_2 \le x_2, \ldots, X_d \le x_d) \\
&= \mu_{\mathbf{X}}\Big( (-\infty,x_1] \times (-\infty,x_2] \times \cdots \times (-\infty,x_d] \Big)
\end{align*}

\end{defn}

\begin{thm}[Properties of distribution functions]
\label{thm-multidim}
If $F=F_{\mathbf X}$ is a distribution function of a
$d$-dimensional random vector, then it has the
following properties:

\begin{enumerate}[label=(\roman*)]
\item $F$ is nondecreasing in each coordinate.

\item For any $1\le i\le d$,
$\lim_{x_i\to-\infty} F(\mathbf{x}) = 0$.

\item $\lim_{\mathbf{x}\to(\infty,\ldots,\infty)}
F(\mathbf{x}) = 1$.

\item $F$ is right-continuous, i.e.,
$F(\mathbf{x}+) := \lim_{\mathbf{y}\downarrow
\mathbf{x}} F(\mathbf{y}) = F(\mathbf{x}),$
where here $\mathbf{y}\downarrow \mathbf{x}$ means
that $y_i\downarrow x_i$ in each coordinate.

\item For $1\le i\le d$ and $a<b$, denote by
$\Delta_{a,b}^{x}$ the \emph{differencing operator in the
variable $x$}, which takes a function $f$ of the real
variable $x$ (and possibly also dependent on other
variables) and returns the value

\[
\Delta_{a,b}^{x}f = f(b) - f(a)
\]

\noindent Then, for any real numbers $a_1 < b_1$, $a_2<b_2$,
\ldots, $a_d < b_d$, we have

\[
\Delta^{x_1}_{a_1,b_1} \Delta^{x_2}_{a_d,b_d}
\ldots \Delta^{x_d}_{a_d,b_d} F \ge 0.
\]
\end{enumerate}

\end{thm}

\begin{proof}
See Chapter 1 in [Dur2010] or Appendix A.2 in
[Dur2004].
\end{proof}

\begin{thm}
Any function $F$ satisfying the properties in \texorpdfstring{\hyperref[thm-multidim]{Theorem~\ref*{thm-multidim}}}{Theorem} above is a distribution function of
some random vector $\mathbf{X}$.
\end{thm}

\begin{proof}
See Chapter 1 in [Dur2010] or Appendix A.2 in
[Dur2004].
\end{proof}
